\section{Rational certificates and compressed feasible trajectories}
\label{sec:dynamics-bit}

For polynomial dynamics, the exact local optimization problems in
\cref{sec:dynamics} can be replaced by rational linear programs.  Exact
forward simulation requires a different representation: polynomial
composition can produce rational states with exponentially long expanded
denominators.  We retain the rational initial state and controls, and let
the original recurrence define the exact trajectory.  Contraction provides
short rational enclosures for its objective value.  Rounding every
coordinate of the next center then prevents denominator growth across
regridding stages.

\begin{assumption}[Rational polynomial data]
\label{ass:dynamics-bit}
Assume \cref{ass:dynamics}.  Every interval endpoint and every coefficient
of $\varphi_t$ and $a_t$ is rational.  Their numerical degrees are bounded
by a fixed constant.  The local dimensions are fixed, with
$N=T$, $p=3$, $k=2$, and maximum separator dimension $q\le1$
($q=0$ when $T=1$).
Valid rational bounds $a,b,H,M$, with $0\le a<1$, are supplied for the
derivatives in \cref{ass:dynamics}.  The box-invariance assertion and
feasible quadratic growth are promises.  The positive growth constant
$g$ need not be supplied.
\end{assumption}

Choose a rational bound $G_b\ge0$ on the absolute value of every component
of every bag-cost gradient on its full box.  At fixed degree this can be
done in polynomial bit work: bound each differentiated monomial by its
coefficient magnitude times the appropriate powers of the maximum
coordinate magnitudes, and sum these bounds.  Set $G=2G_b$.  Since each
state occurs in at most two bags,
\[
 \abs{\partial_{\sigma_j}F(c)}\le G,\qquad j=1,\ldots,T,
\]
for every ambient-box center $c$, whether or not it satisfies the dynamics.
Thus \cref{lem:dynamics-adjoint,lem:dynamics-error} apply to the rational
centers used below.  The bounds $M,H,a,G_b$ and box invariance determine
certificate soundness.  Growth is used only to prove a rate.

\subsection{Exact rational local optimization}

\begin{lemma}[Affine models and exact local linear programs]
\label{lem:bit-local-lp}
For a bag leaf $B$ of width $w_B$, let $v_B$ be its full midpoint and use
\begin{equation}
 m_{t,B}(z)=a_t(v_B)+\inner{\nabla a_t(v_B)}{z-v_B}
                         -\frac{Mp}{8}w_B^2.
 \label{eq:bit-affine-model}
\end{equation}
This rational affine model satisfies \cref{eq:model-contract} with
$A=pM$.  Use the rational graph strip of \cref{lem:dynamics-strip}.
Every finite local task in the certificate dynamic program is a rational
linear program with at most three variables and eight inequalities.  Its
minimum and an attaining point, or its emptiness, can be found exactly by
a constant number of rational basis computations.
\end{lemma}

\begin{proof}
Taylor's estimate on the box gives
\[
 \abs{a_t(z)-a_t(v_B)-\inner{\nabla a_t(v_B)}{z-v_B}}
 \le\frac M2\norm{z-v_B}^2\le\frac{Mp}{8}w_B^2.
\]
Subtracting the last bound proves
$0\le a_t(z)-m_{t,B}(z)\le Mp w_B^2/4$.  No vertex exactness is needed.
The graph strip has the form
\[
 \abs{\sigma'-\varphi_t(v)
      -\inner{\nabla\varphi_t(v)}{(\sigma,u)-v}}
 \le \frac H4w_B^2,
\]
where $v$ is the input midpoint.  All its coefficients are rational.
It contains every true graph point in $B$, and every point of the strip
has residual at most $Hw_B^2/2$.

Intersecting $B$ with its own separator cell only changes coordinate
bounds.  Merge these bounds by taking the maximum lower endpoint and
minimum upper endpoint.  There are then at most six coordinate
inequalities and two strip inequalities.  The model and separator terms
are affine.  Child-message intercepts are constants, so they can be
added after minimizing the nonconstant part.

For an unreduced three-dimensional task, enumerate the at most
$\binom83=56$ triples of constraint rows.  Discard dependent triples,
solve each remaining rational linear system, and check all inequalities
exactly.  Select a feasible candidate of least objective value.  A
nonempty bounded polyhedron has a vertex.  At a vertex its active normals
span the ambient space: otherwise a nonzero common null direction and
its negative give feasible small displacements, contradicting
extremality.  Thus a vertex appears in this enumeration, and finding no
feasible candidate proves emptiness.  This argument also covers
lower-dimensional polyhedra.  Alternatively, eliminate coordinates fixed
by the intersected bounds and enumerate bases in the remaining dimension;
in dimension zero, check the unique point directly.  Touching faces and
the exact affine graph when $H=0$ therefore require no numerical tolerance.

The infeasibility flags and bottom-up rules in \cref{sec:dynamics} remain
valid.  A true feasible trajectory supplies an available configuration,
so a finite root minimum exists.  Exact backtracking returns rational
bag points and hence a rational top-copy point, rational controls, and a
rational initial state.
\end{proof}

\subsection{Enclosing an exact feasible trajectory}

Fix the rational initial state and controls obtained by backtracking.
The equations
$\sigma_{t+1}=\varphi_t(\sigma_t,u_t)$ define the exact states without
expanding their numerators and denominators.  Box invariance makes this
trajectory feasible.

\begin{lemma}[Contraction enclosures and an objective upper bound]
\label{lem:bit-enclosure}
For any dyadic mesh $\delta>0$, rational state intervals can be computed
with midpoints $\widehat\sigma_t$ and radii $r_t$ such that they contain
the exact states, lie in the prescribed state intervals, and satisfy
\begin{equation}
 r_0=0,\qquad r_{t+1}\le ar_t+\delta,
 \qquad r_t\le\frac{\delta}{1-a}.
 \label{eq:bit-enclosure-radius}
\end{equation}
The rational number
\begin{equation}
 U_{\mathrm{enc}}=\sum_{t=0}^{T-1}a_t(\widehat\sigma_t,u_t,\widehat\sigma_{t+1})
               +G_b\sum_{t=0}^{T-1}(r_t+r_{t+1})
 \label{eq:bit-objective-upper}
\end{equation}
satisfies
\[
 F(y)\le U_{\mathrm{enc}}\le F(y)+\frac{4G_bN\delta}{1-a}.
\]
\end{lemma}

\begin{proof}
Start from the singleton interval containing the exact rational initial
state.  Given its current midpoint and radius, evaluate
$\varphi_t(\widehat\sigma_t,u_t)$ exactly as a rational number.  The
derivative bound gives the enclosure
\[
 [\varphi_t(\widehat\sigma_t,u_t)-ar_t,
   \varphi_t(\widehat\sigma_t,u_t)+ar_t].
\]
Round its lower endpoint down and its upper endpoint up to the dyadic
mesh, and intersect it with the prescribed next-state interval.  Box
invariance preserves containment under this intersection.  Outward
rounding increases the radius by at most $\delta$, and intersection
cannot increase it.  This proves the recurrence and its geometric-sum
bound.  The intersection also keeps every subsequent midpoint in the
box where the derivative promise holds.  Ordinary interval substitution
in a polynomial does not establish this width recurrence: repeated
occurrences of the state can overestimate variation.

Both the exact bag point and its midpoint approximation lie in the full
bag box.  Integrating the componentwise gradient bound along their segment
gives
\[
 \abs{a_t(\widehat\sigma_t,u_t,\widehat\sigma_{t+1})-a_t(y_{V_t})}
 \le G_b(r_t+r_{t+1}).
\]
Adding these inequalities proves that $U_{\mathrm{enc}}$ is an upper bound and that
its excess is at most $2G_b\sum_t(r_t+r_{t+1})$, hence at most the stated
$4G_bN\delta/(1-a)$.  No independence between adjacent state intervals
is required.
\end{proof}

Fix a positive rational $\omega$; one may take $\omega=1$ throughout.
At stage $j$, choose the largest reciprocal power-of-two mesh satisfying
\begin{equation}
 \delta_j\le
 \min\left\{1,
       \frac{\omega(1-a)h_j^2}{4\max\{1,G_b\}},
       \frac{(1-a)h_{j+1}}4\right\}.
 \label{eq:bit-enclosure-mesh}
\end{equation}
The computed upper bound $U_j$ then exceeds $F(y_j)$ by at most
$\omega Nh_j^2$.  Keep the least such upper bound, together with the
rational initial state and controls that produced it.

\subsection{Rational center resetting}

The feasible incumbent and the next center serve different purposes.
The incumbent retains its exact rational initial state and controls.
The center retains only newly rounded coordinates and may be infeasible.

\begin{lemma}[Resetting and the rounded-center rate]
\label{lem:bit-reset}
At stage $j$, use the rational center $c_j$, the models of
\cref{lem:bit-local-lp}, and the exact adjoints of
\cref{lem:dynamics-adjoint}.  Let $y_j$ be the forward repair of the
backtracked top-copy point.  Round its enclosure midpoints and its exact
input coordinates to the largest reciprocal power-of-two mesh of size at
most $\min\{1,h_{j+1}/2\}$, then clip
each rounded coordinate to its exact interval.  The resulting center
satisfies
\begin{equation}
 \norm{c_{j+1}-y_j}^2\le pNh_{j+1}^2.
 \label{eq:bit-reset}
\end{equation}
Suppose $\theta$ satisfies \cref{eq:dynamics-grading}, with $G=2G_b$
and $A=pM$.  Write $C$ for the corresponding constant in
\cref{eq:dynamics-constants}, and put
\begin{equation}
 B_{\mathrm{round}}=\max\left\{p,
                 \frac{4(C+\omega)}g+\frac{3p}5\right\}.
 \label{eq:bit-rounded-constant}
\end{equation}
Starting from the rational box midpoint, the successive repaired points
obey $e_j=\norm{y_j-x^*}^2\le B_{\mathrm{round}}Nh_j^2$ and
\begin{equation}
 0\le\UB-\LB_j\le gB_{\mathrm{round}}Nh_j^2.
 \label{eq:bit-gap}
\end{equation}
\end{lemma}

\begin{proof}
By \cref{eq:bit-enclosure-mesh}, each state midpoint has error at most
$h_{j+1}/4$.  Rounding to the prescribed mesh adds at most
$h_{j+1}/2$.  The initial state and controls have only the latter error.
Clipping cannot increase distance from a coordinate in the interval.
Every coordinate error is therefore at most $h_{j+1}$.  There are
$2T+1\le pN$ coordinates, proving \cref{eq:bit-reset}.

For initialization, let $c_0$ be the rational box midpoint and choose any
feasible reference trajectory $y_{-1}$, for example the recurrence from
its retained input coordinates.  Both $c_0$ and $y_{-1}$ lie in the box,
so $\norm{c_0-y_{-1}}^2\le pNh_0^2$ and
$e_{-1}\le pNh_0^2$.  This reference is needed only for the proof;
its expanded coordinates and objective value are not computed.

The arbitrary-center estimate \cref{eq:dynamics-error-final} gives
\[
 F(y_j)-\LB_j\le\frac g{20}\norm{y_j-c_j}^2+CNh_j^2.
\]
The three-vector inequality and the reset bound give
$\norm{y_j-c_j}^2\le3(e_j+e_{j-1}+pNh_j^2)$, including stage zero
by the initialization bound.  Since $\LB_j\le f^*$, feasible quadratic
growth yields
\begin{equation}
 e_j\le\frac3{17}e_{j-1}
              +\frac{20C/g+3p}{17}Nh_j^2.
 \label{eq:bit-rounded-recursion}
\end{equation}
In particular the contraction factor $3/17$ is smaller than $1/4$.
For $j\ge1$, induction gives
$e_{j-1}\le4B_{\mathrm{round}}Nh_j^2$; this bound also holds at
$j=0$ because $B_{\mathrm{round}}\ge p$.  The induction closes because
\[
 12B_{\mathrm{round}}+20C/g+3p\le17B_{\mathrm{round}}.
\]
Adding the upper-bound evaluation error and using these distance bounds
gives
\begin{align*}
 0\le\UB-\LB_j
 &\le\left(\frac{3gB_{\mathrm{round}}}4+\frac{3gp}{20}
                        +C+\omega\right)Nh_j^2\\
 &\le gB_{\mathrm{round}}Nh_j^2,
\end{align*}
where the last inequality follows from
$gB_{\mathrm{round}}/4\ge C+\omega+3gp/20$.
\end{proof}

The center reset includes the controls and initial state.  Retaining their
arbitrary LP denominators in successive centers would defeat the bit-length
argument even if all dependent states were rounded.  At an infeasible
center we still use the original cost models and the adjusted gradients
only for separator slopes.  A model of the full adjusted function would
also need its nonzero constant $\sum_t\mu_tq_t(c)$; the present construction
uses the exact identity in \cref{lem:dynamics-error} directly.

\subsection{Bit lengths and conditioned complexity}

Let $I$ be the total binary length of the rational input, including supplied
bounds and tolerance.  The fixed rational $\omega$ is included when it is
not set to one.  Rational numbers are stored as reduced numerator-denominator
pairs.  Write $\theta=2^{-\mu}$, $\mu\ge1$.

\begin{lemma}[Polynomial numerical representations]
\label{lem:bit-size}
At stage $j$, let $b_j=O(I+j+\mu+\log(T+1))$.  Reset-center coordinates
have $O(I+j)$ bits, fresh shell endpoints have $O(I+j+\mu)$ bits, adjoints
and local affine contributions have $O(Tb_j)$ bits, and every finite
dynamic-programming message has $O(T^2b_j)$ bits.  Forward enclosures and
their objective upper bounds also have polynomial bit length.  Each local
LP and each arithmetic operation can therefore be performed in polynomial
bit work in $T,I,j,\mu$.
\end{lemma}

\begin{proof}
A reset coordinate is a dyadic rational at the prescribed mesh or a clipped
input endpoint; its magnitude is bounded by the input box.  The required
mesh exponent is $O(I+j)$.  Fresh shell boundaries are formed from these
centers, $h_j=s_0 2^{-j}$, and the dyadic grading ratio, so their bit
lengths are $O(I+j+\mu)$.  They do not inherit denominators from the
previous stage's LP solutions.

At fixed degree and arity, rational polynomial values and derivatives at
points of this size have $O(b_j)$-bit representations.  In the adjoint
recurrence
$\mu_{t-1}=A_t\mu_t-d_t$, each coefficient $A_t,d_t$ has $O(b_j)$
bits.  Multiplication by $A_t$ and addition of $d_t$ increase length by
only $O(b_j)$ per step.  Across $T$ steps this gives $O(Tb_j)$ bits.
Fixed-size vertex determinants and local objective evaluation consequently
have polynomial bit length; the bound $O(Tb_j)$ suffices for both local
points and nonconstant local values.

A finite message is a sum of chosen local contributions over its subtree.
To see this directly, write a node's contribution as its affine model plus
its child slopes evaluated at the node point, minus its own separator
slope evaluated there.  The dynamic-programming recurrence adds selected
child intercepts.  Induction expresses the resulting message as a sum
with one such contribution per selected bag, hence at most $T$ terms.
Taking a minimum only selects one of these sums.  No two messages are
multiplied.  A sum of $T$ rational numbers of length $O(Tb_j)$ has length
$O(T^2b_j)$ after reduction, regardless of how many grid tasks were
compared.  Exact comparisons and rational reduction take polynomial work
at these lengths.

One can choose \cref{eq:bit-enclosure-mesh} with
\[
 \log(1/\delta_j)=
 O\left(I+j+\log(1+G_b)+\log\frac1{1-a}
                         +\log(1+1/\omega)\right).
\]
The scale terms have input-sized binary encodings.  At each forward step,
polynomial evaluation uses the current bounded-length midpoint and fixed
rational control.  Outward rounding then returns to the prescribed dyadic
precision; clipping can introduce only input endpoint denominators.
There is no nonlinear composition of unrounded rational state values.
The initial-state singleton and controls have the polynomial lengths
already proved for LP points.  Finally, summing the $T$ rational cost
values and radii in \cref{eq:bit-objective-upper} still gives polynomial
length.
\end{proof}

If a rational growth lower bound is supplied, all grading choices can be
made by rational arithmetic.  For $p=3$, $k=2$, use the rational majorants
\[
 \widehat L_q=1+a+b,\qquad
 \widehat\rho=\frac{3(1+a+b)+(H/2)s_0}{1-a},\qquad
 \widehat D=(3+2\widehat\rho)^2.
\]
They dominate the corresponding $L_q,\rho,D$ in
\cref{eq:dynamics-constants}.  Substituting $\widehat D$ for $D$ in
$B_0,C$ and the grading conditions preserves every estimate.  Square the
nonnegative sides of the condition containing $\sqrt D$, and repeatedly
halve a dyadic ratio until all the resulting rational inequalities hold.
The stopping-stage estimate can likewise be computed by comparing
$gB_{\mathrm{round}}Ns_0^2 4^{-J}$ with $\eps$.

\begin{theorem}[Conditioned rational approximation]
\label{thm:dynamics-bit}
Under \cref{ass:dynamics-bit}, fix a sufficient dyadic ratio and run the
rational local optimization, enclosure, and center-reset procedure above.
For every rational $\eps>0$, it returns a finite rational lower certificate,
a rational feasible-objective upper bound, and rational initial state and
controls in their exact intervals, with $\UB-\LB\le\eps$.
The original recurrence is part of the output and defines the exact feasible
trajectory from these retained inputs.  No expanded state numerators or
denominators are required.

The algorithm stops by
\[
 J=\max\left\{0,
       \ceil{\log_2\left(s_0\sqrt{gB_{\mathrm{round}}N/\eps}\right)}
       \right\},
\]
has the sharp shell counts of \cref{eq:dynamics-counts} with this $J$,
and uses at most
\begin{equation}
 N(4/\theta)^3(J+1)^3\operatorname{poly}(T,I,J,\mu)
 \label{eq:bit-work}
\end{equation}
bit operations.  Choosing the first sufficient dyadic ratio, or the first
one satisfying the conservative rational tests above, makes this work
polynomial in input size, horizon, $\log_+(1/\eps)$, and the numerical
conditioning parameters.  One may take $0<g\le1$ without loss and,
with $\omega=1$,
use
\[
 \mathcal K=1+b+H+G_b+M+s_0+\frac1{1-a}+\frac1g
\]
as a single conditioning bound.  The assertion is polynomial in these
numerical quantities, rather than uniformly polynomial in their binary
encoding lengths.
\end{theorem}

\begin{proof}
Each stage has a valid lower certificate by \cref{sec:dynamics} and
\cref{lem:bit-local-lp}.  Its compressed repaired trajectory is feasible
by box invariance, and its upper bound is sound by
\cref{lem:bit-enclosure}.  Thus the rational stopping test proves the
claimed objective gap.  \Cref{lem:bit-reset} gives the stated stage bound.
The incumbent stores the retained input coordinates associated with its
least upper bound, so its trajectory also satisfies the returned gap.

The fresh partitions and incidence lists have the counts established in
\cref{prop:regridding-incidence,eq:dynamics-counts}.  Here the actual
separator dimension is at most one.  Empty domains only remove available choices,
and each local task now requires a constant number of rational basis solves.
\Cref{lem:bit-size} bounds their arithmetic costs.  Grid generation,
gradients, adjoints, forward enclosures, resetting, and output construction
also take polynomial work in the same sizes, proving \cref{eq:bit-work}.

The constants $D,\overline M,B_0,C$ are bounded by polynomials in
$\mathcal K$ after replacing $g$ by $\min\{g,1\}$.  The first sufficient
dyadic reciprocal ratio is also polynomially bounded in $\mathcal K$ by
\cref{eq:dynamics-grading}.  Hence $B_{\mathrm{round}}$ is polynomially
bounded and $J$ grows only logarithmically with the numerical scale, horizon,
and reciprocal tolerance.  This proves the conditioned complexity claim.

The certificate's numerical entries and flags have the bit lengths in
\cref{lem:bit-size}.  A checker can reproduce each constant-size LP basis
enumeration, verify finite message recurrences and empty-state flags, and
run the upper-bound enclosure recurrence.  It trusts, or separately
verifies, the supplied global derivative bounds and box-invariance promise.
Thus proof serialization and these local checks fit the polynomial work
bound. For any further rational requested state precision $\tau>0$, rerunning
\cref{lem:bit-enclosure} with $\delta\le(1-a)\tau$ encloses every
retained trajectory state with radius at most $\tau$, in bit work
polynomial in the preceding parameters and the requested precision length.
\end{proof}

\subsection{Searching without a supplied growth constant}

\begin{theorem}[Unknown-growth search with bit budgets]
\label{thm:dynamics-unknown-growth}
Under \cref{ass:dynamics-bit}, there is a sound algorithm that uses no
supplied value of $g$ and has polynomial bit complexity in the parameters
of \cref{thm:dynamics-bit}.  Its returned certificate and compressed
feasible trajectory have polynomial output length.  The sharp shell-count
conclusion of \cref{thm:dynamics-bit} concerns the sufficient fixed-ratio
run; the search algorithm's output bound follows from its total bit work.
\end{theorem}

\begin{proof}
Trial $\mu\ge1$ uses $\theta=2^{-\mu}$ and starts at the rational box
midpoint.  It executes stages $j=0,1,\ldots$ using the rational LPs,
exact adjoints, enclosure mesh, and reset procedure above.  These operations
use only the input bounds, $h_j$, and the fixed $\omega$.  They use none
of $g,C,B_{\mathrm{round}}$, or $J$.  A trial succeeds only when its
computed rational upper and lower bounds have gap at most $\eps$.
Every such success is sound, whether or not the ratio satisfies the
convergence inequalities.

In round $r=1,2,\ldots$, restart each trial $1\le\mu\le r$ and simulate
at most $2^r$ elementary bit operations of its algorithm.  Include setup,
grid generation, arithmetic, control logic, and completed output in this
trial budget; no trial-specific preprocessing occurs outside it.  Stop an
unsuccessful simulation when its budget expires, even within an arithmetic
routine.  A bit-operation budget imposes no separate stage cap and therefore
also handles trials whose late stages are expensive.  A binary clock and
time-bounded simulation add at most a polynomial logarithmic overhead to
the simulated work.  Completed output is part of the successful trial's
work.

Let $\mu^*$ be the smallest sufficient dyadic index.  With
$\eta=g/20$, it satisfies $2^{\mu^*}\le2R$, where
\[
 R=\max\left\{2,\sqrt{12D},
       \frac{4\sqrt{12}\,k\sqrt D\,\overline M}{\eta},
       \sqrt{\frac{48kB_0}{\eta}}\right\}.
\]
Indeed, every reciprocal ratio at least $R$ satisfies each condition in
\cref{eq:dynamics-grading}, and the next power of two is at most $2R$.
Thus $2^{\mu^*}$ is polynomially bounded in the conditioning parameters.
Let $W^*\ge1$ bound the elementary bit work of this trial through its
successful output, as supplied by \cref{thm:dynamics-bit}.  Set
\[
 r^*=\max\{\mu^*,\ceil{\log_2W^*}\}.
\]
Round $r^*$ contains the sufficient trial and supplies enough work for it
to succeed, unless another sound trial has already returned.  Moreover,
\begin{equation}
 \sum_{r=1}^{r^*}r2^r\le2r^*2^{r^*}
       \le4r^*\max\{2^{\mu^*},W^*\}.
 \label{eq:bit-budget}
\end{equation}
Since $r^*$ is logarithmic in the maximum on the right, accounting for the
clock and simulation overhead preserves polynomial bit complexity.
Writing the returned output was budgeted, so its length is bounded by the
same total work.

A coarse-ratio trial can succeed before the sufficient trial and may do
so at a different stage.  Its stage index has not been bounded by the
sufficient trial's $J$.  Consequently the search gives a polynomial
output bound, while the sharper final leaf-and-cell count applies to the
fixed sufficient ratio.  No growth test or estimate of the best growth
constant is needed.
\end{proof}

This result uses rational polynomial data of fixed numerical degree and
the compressed feasible-solution representation.  It does not provide an
expanded rational coordinate vector of polynomial length.  For example,
$\varphi(\sigma)=\sigma^2/4$ on $[0,1]$, starting at $\sigma_0=1/2$,
has derivative bounded by $1/2$ but gives
$\sigma_t=2^{-(3\cdot2^t-2)}$.  Its reduced denominator is exponentially
long, whereas its recurrence and rational initial state are short.
Additional constraints require a repair that preserves them; the constraint
obstruction in \cref{prop:limitations-constraint} lies outside the present
box-invariant model.
